\subsection{R-32: unpadded common stability and nonrectangular period one}

\paragraph{Statement R-32.}
Let $d\geq3$, $4\leq r\leq2d-2$, and let
$\nu\vdash ds$, $s\geq1$, have exactly $r$ positive parts.
Put $\eta=\nu-\nu_r\mathbf1_r$. The enhanced cones for the
\emph{unpadded} degree-$d$ permanent and determinant images
in $r$ variables have a common stable point. In particular,
their enhanced invariant dimensions equal their actual
image dimensions plus $r_\eta-(r^2-1)$, and the exact
single-label extraction in G-22 applies.

If $\nu$ is nonrectangular, these two cones also have a common
stable point $w$ satisfying
\[
 gw=cw,\quad g\in\mathrm{SL}_r,\quad c\in\mathbb C^*
             \quad\Longrightarrow\quad c=1.
\]
Consequently both unpadded multiplicities at $j\nu$, in
coordinate degree $sj$, are positive at every integer
$j\geq\beta^2$ for any common ambient invariant-generator
bound $\beta$. The rectangular case retains the stable-point
and exact-dimension assertion; no period-one conclusion is
asserted there.

\begin{proof}
Take $V=\mathbb C^r$, with its standard Hermitian metric,
and set $\mu_i=\nu_i/s$. These are positive rational
numbers with sum $d$. Split $V=V_1\oplus V_2$ orthogonally
into blocks of sizes $\lfloor r/2\rfloor$ and
$\lceil r/2\rceil$. Each block has dimension at least two
and at most $d-1$.

Let $w_b$ be the sum of the $\mu_i$ in block $b$, divided
by $d$. Then $w_1,w_2>0$ and $w_1+w_2=1$. The prescribed
product-moment lemma in R-29, with its polynomial degree
there renamed $d$, supplies unit degree-$d$ products $P,Q$
in the respective blocks such that their moment matrices
are the corresponding diagonal lists $\mu_i/w_b$.
Both products have finite complex projective stabilizer
and independent constant directional derivatives.
For
\[
                   q=\sqrt{w_1}P+\sqrt{w_2}Q,
\]
the orthogonal-block calculation in R-29 gives
\[
                  M(q)=\operatorname{diag}(\mu_1,\ldots,\mu_r).
\]
Its complex projective stabilizer is finite as well. Indeed,
the two mixed infinitesimal bidegrees $(d-1,1)$ and
$(1,d-1)$ are distinct for $d\geq3$. Derivative independence
kills both off-diagonal Lie blocks. Finite projective
stabilizers of $P,Q$ make the two remaining blocks scalar,
and the equal polynomial degrees force one common scalar.
The two-product cycle construction of R-29 represents this
same $q$ as a size-$d$ permanent and as a size-$d$ determinant,
with the appropriate cycle sign.

For the standard unit auxiliary extremal vector $v_\eta$,
consider
\[
                    w_0=q^{\otimes s}\otimes v_\eta.
\]
Its full moment is
\[
 sM(q)-\operatorname{diag}(\eta)=\nu_r I_r.
\]
A projective symmetry of this pure tensor preserves $[q]$,
so its projective stabilizer is finite. The cited
Kempf--Ness criterion (G-B14) therefore makes $w_0$ stable.
No padding moment or first-row inequality is involved.
This proof also works when $\eta=0$, proving the stable
assertion for rectangular $\nu$.

Suppose now that $\eta\ne0$, and fix this $q$. Let $H$ be
its finite projective automorphism group. The auxiliary
extremal lines form a nontrivial irreducible partial flag
orbit, on which $\mathrm{PGL}_r$ acts faithfully, as checked
in R-30. Stability is open, and the fixed-$q$ stable flag
locus contains the flag of $w_0$. Each nonidentity element
of $H$ fixes a proper closed subset of the flag orbit.
Choose a stable flag outside their finite union.

If $gw=cw$ for the resulting tensor, the projective class
of $g$ preserves both $[q]$ and this flag and is therefore
the identity. Hence $g$ is central in $\mathrm{SL}_r$.
Its exponent on the enhanced tensor is
\[
                       ds-|\eta|=r\nu_r,
\]
so the center acts trivially and $c=1$. The closed-orbit
separation and numerical-semigroup argument in R-28 now
gives the asserted all-integer occurrence bound. Finally,
the exact coordinate extraction and stable quotient-dimension
calculation of G-22 do not require padding and apply to the
two unpadded cones just used.
\end{proof}

\paragraph{Scope.}
This lemma supplies stable unpadded comparison algebras even
when an available smaller padded form fails its first-row
strict inequality. Its rectangular part supplies a stable
quotient of the stated dimension, while allowing a nontrivial
invariant-degree period. It gives no equality between the two
multiplicity functions and does not extend the row range
beyond $2d-2$.
